\begin{abstract}
5G virtualized radio access networks (vRAN) increasingly run physical-layer (PHY) processing on GPUs to meet sub-millisecond slot deadlines. Current deployments consolidate multiple cells onto a single GPU, expecting that compute parallelism will absorb the load. We show that this assumption is wrong. Through systematic profiling of the NVIDIA Aerial cuPHY uplink pipeline on an A100 GPU, we find that the dominant kernel---channel equalization coefficient application---is memory-bound at an arithmetic intensity of 4.32~FLOP/byte, achieving only 0.86\% of peak Tensor Core throughput even after FP16 optimization. The bottleneck is not compute but memory bandwidth: L1 cache throughput saturates at 93.6\% in the scalar baseline and DRAM throughput rises to 21.2\% after Tensor Core acceleration, leaving no headroom for concurrent cells.

When multiple cells share a GPU via CUDA streams, memory bandwidth contention causes per-cell latency degradation of up to 7.7$\times$ at 8~concurrent cells---a superlinear growth that current schedulers, which allocate based on SM occupancy, cannot prevent. Through six controlled experiments on an A100, we establish three key findings: (1)~the contention threshold is the L2~cache capacity (42~MB), not HBM bandwidth---two cells whose combined working set fits in L2 degrade only 1.12$\times$, while cells exceeding L2 degrade 1.50$\times$ or more; (2)~pipeline stage ordering matters---launching the smaller-working-set kernel first reduces degradation from 1.50$\times$ to 1.16$\times$; (3)~online contention prediction from just 5~runtime latency samples achieves $<$3\% error, enabling adaptive batch sizing without offline profiling or hardware PMUs.

Based on these findings, we propose a research line for a memory-hierarchy-aware scheduler combining L2-aware stage co-scheduling, online contention prediction, and adaptive batch sizing with deadline guarantees. We present the experimental evidence, the contention model, and the development plan toward a full system targeting NSDI.
\end{abstract}
